\documentclass[11pt]{article}

\usepackage{multicol}

\usepackage[french, english]{babel}

\usepackage[T1]{fontenc}

\usepackage{amsfonts}
\usepackage{amsmath}
\usepackage{amssymb}

\usepackage{graphicx}
\usepackage{multirow}
\usepackage{multicol}
\usepackage{xcolor}
\usepackage[shortlabels]{enumitem}

\usepackage[margin=2cm]{geometry}

\usepackage{mathtools}
\usepackage{siunitx}
\usepackage{braket}

\usepackage{listings}

\usepackage{tikz}
\usetikzlibrary{shapes.geometric, arrows.meta, positioning}

\tikzstyle{block} = [rectangle, rounded corners, minimum width=4cm, minimum height=1cm,
text centered, draw=black, fill=blue!10]
\tikzstyle{process} = [rectangle, rounded corners, minimum width=5cm, minimum height=1cm,
text centered, draw=black, fill=green!10]
\tikzstyle{io} = [trapezium, trapezium left angle=70, trapezium right angle=110,
minimum width=4cm, minimum height=1cm, text centered, draw=black, fill=orange!15]
\tikzstyle{decision} = [diamond, aspect=2, draw=black, fill=yellow!15, text centered]
\tikzstyle{arrow} = [thick, ->, >=Stealth]

\renewcommand\lstlistingname{Code Python}

\lstset{frame=tb,
  language=Python,
  numbers=left,
  stepnumber=1,
  aboveskip=3mm,
  belowskip=3mm,
  showstringspaces=false,
  columns=flexible,
  basicstyle={\small\ttfamily},
%   numbers=none,
  numberstyle=\tiny\color{gray},
  keywordstyle=\color{blue},
  commentstyle=\color{paired2b},
  stringstyle=\color{paired3b},
  breaklines=true,
  breakatwhitespace=true,
  tabsize=1,
  frame=single,
%   showtabs=false,
}


\begin{document}

%\renewcommand*{\contentsname}{Table des matières}

%\hrule

%\begin{center}
  %{\Large STT290 - Probabilitées} \\ \vspace{5mm}
  %{\LARGE\sffamily Résumé} \\
  %Alexis Jean\\
%\end{center}

%\hrule

%\tableofcontents

\begin{multicols}{2}
  \section{\underline{Variables aléatoires discrètes}}
  \textbf{Variable de Bernoulli :}
  On défini X comme la variable aléatoire où :\\
  $X = \begin{cases}
  1 & \text{avec probabilité p}\\
  0 & \text{avec probabilité 1-p}
  \end{cases}$\\
  $\mathbb{E}(X)=1\cdot p+0\cdot (1-p)=p$\\
  $\mathbb{V}ar(X)=1^2\cdot p+0^2\cdot (1-p)-p^2=p-p^2=p(1-p)$\\

  \textbf{Variable binomiale :}
  La variable X possède n épreuves avec p probabilité de succès qu'on note :\\
  $X~Bin(n,p)$\\
  $\mathbb{P}(X=k)=\binom{n}{k}p^k(1-p)^{n-k}, k=0,1, ... ,n$\\
  $F(k)=\mathbb{P}(X\leq k)=\sum_{i\leq k}\binom{n}{i}p^i(1-p)^{n-i}$, \\où $i=0,1, ... ,n$\\
  $\mathbb{E}(X)=np$, $\mathbb{V}ar(X)=np(1-p)$\\

  \section{\underline{Variables aléatoires continues}}


  \textbf{Variable aléatoire continues :} Une variable aléatoire $X$ est dite continue si elle peut prendre une infinité de valeurs dans un intervalle de $\mathbb{R}$. \\
  $\mathbb{P}(x \epsilon B) = \int_{x \epsilon B} f(x) dx$ \\
  

  \textbf{Fonction de répartition :} La fonction de répartition $F(x)$ d'une variable aléatoire continue $X$ est définie par : \\
  $F(x) = \mathbb{P}(X \leq x) = \int_{-\infty}^{x} f(t) dt$ \\
  Donc, la fonction de densité $f(x)$ est la dérivée de la fonction de répartition $F(x)$ : \\
  $f(x) = \frac{d}{dx} F(x)$ \\
  
  $\mathbb{E}(X) = \int_{-\infty}^{\infty} x f(x) dx$ \\
  $Var(X) = \int_{-\infty}^{\infty} (x - E(X))^2 f(x) dx$ \\


  \textbf{Variable aléatoire uniforme continue :} Une variable aléatoire $X$ est dite uniforme continue sur l'intervalle $[a, b]$ si sa fonction de densité est donnée par : \\
  \[
  f(x) =
  \begin{cases}
    \dfrac{1}{b - a} & \text{si } a \leq x \leq b, \\
    0 & \text{sinon}
  \end{cases}
  \]
  La fonction de répartition est : \\
  \[
  F(x) =
  \begin{cases}
    0 & \text{si } x < a, \\
    \dfrac{x - a}{b - a} & \text{si } a \leq x \leq b, \\
    1 & \text{si } x > b
  \end{cases}
  \] \\


  \textbf{Variable aléatoire Beta :} Une variable aléatoire $X$ suit la loi Beta si sa fonction de densité est donnée par : \\
  \[
  f(x) =
  \begin{cases}
    \dfrac{1}{B(a, b)} x^{a-1} (1-x)^{b-1} & \text{si } 0 \leq x \leq 1, \\
    0 & \text{sinon}
  \end{cases}
  \]
  où les paramètres $a, b > 0$ et $B(a, b)$ = $\int_{0}^{1} x^{a-1} (1-x)^{b-1} dx$ est la fonction bêta. \\

  $\mathbb{E}(X) = \dfrac{a}{a + b}$ \\
  $Var(X) = \dfrac{ab}{(a + b)^2 (a + b + 1)}$ \\


  \textbf{Variable exponentielle :} Une variable aléatoire $X$ suit la loi exponentielle avec le paramètre $\lambda > 0$ si sa fonction de densité est donnée par : \\
  \[
  f(x) =
  \begin{cases}
    \lambda e^{-\lambda x} & \text{si } x \geq 0, \\
    0 & \text{sinon}
  \end{cases}
  \]

  La fonction de répartition est : \\
  \[
  F(x) =
  \begin{cases}
    0 & \text{si } x < 0, \\
    1 - e^{-\lambda x} & \text{si } x \geq 0
  \end{cases}
  \] 

  $\mathbb{E}(X) = \dfrac{1}{\lambda}$,
  $Var(X) = \dfrac{1}{\lambda^2}$\\

  
  \textbf{Variable gamma :} Une variable aléatoire $X$ suit la loi gamma avec les paramètres $\gamma$ et $\lambda$ si sa fonction de densité est donnée par : \\
  \[
  f(x) =
  \begin{cases}
    \dfrac{{\lambda e^{-\lambda x}(\lambda x)^{\gamma - 1}}} {\Gamma(\gamma)} & \text{si } x \geq 0, \\
    0 & \text{sinon}
  \end{cases}
  \]
  où $\Gamma(\gamma) = \int_{0}^{\infty} x^{\gamma - 1} e^{-x} dx$ est la fonction gamma. \\
  Truc : $\Gamma(n) = (n-1)!$ si n est un entier positif. \\
  La fonction de répartition est : \\
  \[
  F(x) =
  \begin{cases}
    \int_{0}^{x} \dfrac{\lambda e^{-\lambda t}(\lambda t)^{\gamma - 1}}{\Gamma(\gamma)} dt & \text{si } x \geq 0, \\
    0 & \text{sinon}
  \end{cases}
  \] 
  $\mathbb{E}(X) = \dfrac{\gamma}{\lambda}$, 
  $Var(X) = \dfrac{\gamma}{\lambda^2}$ \\

  \textbf{Variable de Weibull :} Une variable aléatoire $X$ suit la loi de Weibull si sa fonction de répartition est donnée par : \\
  \[
  F(x) = 
  \begin{cases}
    1 - e^{-(\dfrac{x-\upsilon}{\alpha})^\beta} & \text{si } x \geq \upsilon, \\
    0 & \text{sinon}
  \end{cases}
  \] \\
  où $\alpha, \beta > 0$ et $-\infty < \upsilon < +\infty$. \\ 

  La fonction de densité est donnée par : \\
  \[
  f(x) =
  \begin{cases}
    \dfrac{\beta}{\alpha} \left( \dfrac{x-\upsilon}{\alpha} \right)^{\beta-1} e^{-(\dfrac{x-\upsilon}{\alpha})^\beta} & \text{si } x \geq \upsilon, \\
    0 & \text{sinon}
  \end{cases}
  \]

  $\mathbb{E}(X) = \alpha \Gamma(\dfrac{1}{\beta} + 1)$ \\
  $Var(X) = \alpha^2 \Gamma(\dfrac{2}{\beta} + 1) - \alpha^2 \Gamma^2(\dfrac{1}{\beta} + 1)$ \\


  \textbf{Variable Laplace :} Une variable aléatoire $X$ suit la loi de Laplace si sa fonction de densité est donnée par : \\
  $f(x) =\dfrac{1}{2\sigma} e^{-\dfrac{|x - \mu|}{\sigma}}$ \\
  où $\mu \epsilon \mathbb{R}$ et $\sigma > 0$. \\
  La fonction de répartition est : \\
  \[
  F(x) =
  \begin{cases}
    \dfrac{1}{2} e^{\dfrac{x - \mu}{\sigma}} & \text{si } x < \mu, \\
    1 - \dfrac{1}{2} e^{-\dfrac{x - \mu}{\sigma}} & \text{si } x \geq \mu
  \end{cases}
  \] 
  $\mathbb{E}(X) = \mu$, 
  $Var(X) = 2\sigma^2$ \\


  \textbf{Variable normale (Gaussienne) :} Une variable aléatoire $X$ suit la loi normale si sa fonction de densité est donnée par : \\
  $f(x) = \dfrac{1}{\sigma \sqrt{2\pi}} e^{-\dfrac{1}{2} \left( \dfrac{x - \mu}{\sigma} \right)^2}$ \\
  où $\mu \epsilon \mathbb{R}$ et $\sigma > 0$. On note $X~\sim N(\mu, \sigma^2)$. \\
  La fonction de répartition est : \\
  $F(x) = \dfrac{1}{\sqrt{2\pi}\sigma} \int_{-\infty}^{x} e^{-\dfrac{(t - \mu)^2}{2\sigma^2}} dt$ \\
  $\mathbb{E}(X) = \mu$,
  $Var(X) = \sigma^2$ \\
  (*Voir table de la loi normale ainsi que les théorèmes associés pour les valeurs de $F(x)$*) \\


  \section{\underline{Vecteurs aléatoires}}

  \textbf{Vecteurs aléatoires :} Soient $X$ et $Y$ deux variables aléatoires continues. Le vecteur aléatoire $(X, Y)$ est défini par sa fonction de répartition conjointe $F(x, y)$ telle que : \\
  $F(x,y) = \mathbb{P}(X \leq x, Y \leq y)$ où $x, y \in \mathbb{R}$. \\

  \textbf{Lois conjointes discrètes :} Pour un couple discrèt $(X,Y)$ :\\
  Probabilité conjointe : $p(x,y) = P(X = x, Y = y)$\\
  Marginales : $p_X(x)=\sum_{x}p(x,y), p_Y(y)=\sum_{x}p(x,y)$\\

  \textbf{Densité des vecteurs aléatoire continus :} La fonction de densité conjointe $f(x, y)$ du vecteur aléatoire $(X, Y)$ est définie par : \\
  $\mathbb{P}((X,Y)\in A) = \int\int_{A}f(x,y)dxdy$\\
  Il faut que la densité f respecte les deux conditions suivantes : \\
  1. $f(x,y) \geq 0$ \\
  2. $\int \int_{\mathbb{R}^2} f(x,y) dx dy = 1$ \\
  La relation entre la densié et la distribution est :\\
  $F(x,y)=\mathbb{P}(X \leq a, Y \leq b)=\int_{\infty}^{b}\int_{-\infty}^{a}f(x,y)dxdy$\\
  Donc, $f(x,y)=\frac{\partial^2}{\partial x\partial y}F(x,y)$\\
  Calcul de Probabilité :\\
  Ex.: $\mathbb{P}(x > 1, Y < 1)=2\int_{0}^{1}\int_{1}^{+\infty}e^{-x}e^{-2y}dxdy$\\

  \textbf{Transformations de vecteurs :} Soit $(Y_1, Y_2)=g(X_1, X_2)$. La densité se tranforme via le jacobien :\\
  $f_{Y_1,Y_2}(y_1,y_2)=f_{X1,X2}(h_1(y_1,y_2), h_2(y_1,y_2)) |J_{(h_1,h_2)}|$\\
  Hyp.1: $\begin{cases}
  y_1 = g_1(x_1,x_2)\\
  y_2 = g_2(x_1,x_2)\\
  \end{cases}$ $\implies \begin{cases}
  x_1=h_1(y_1,y_2)\\
  x_2=h_2(y_1,y_2)\\
  \end{cases}$\\
  Hyp.2: \\
  $J_{(g_1,g_2)}(x_1,x_2) = \bigg|\begin{pmatrix}
    \frac{\partial g_1(x_1,x_2)}{\partial x_1} & \frac{\partial g_1(x_1,x_2)}{\partial x_2}\\
    \frac{\partial g_2(x_1,x_2)}{\partial x_1} & \frac{\partial g_2(x_1,x_2)}{\partial x_2}
  \end{pmatrix}\bigg| \neq 0$\\
  Ou encore : \\
  $J_{(h_1,h_2)}(y_1,y_2) = \bigg|\begin{pmatrix}
    \frac{\partial h_1(y_1,y_2)}{\partial y_1} & \frac{\partial h_1(y_1,y_2)}{\partial y_2}\\
    \frac{\partial h_2(y_1,y_2)}{\partial y_1} & \frac{\partial h_2(y_1,y_2)}{\partial y_2}
  \end{pmatrix}\bigg|$\\


  \section{\underline{Quelques propriétés importantes}}

  \textbf{Espérance d'une somme :}\\
  Idée générale : $\mathbb{E}(X+Y)=\mathbb{E}(X)+\mathbb{E}(Y)$\\
  Dans le cas discret : \\$\mathbb{E}(g(X,Y)) = \sum_{x}\sum_{y}g(x,y)p(x,y)$\\
  Dans le cas continu : \\$\mathbb{E}(g(X,Y))=\int\int g(x,y)f(x,y)dxdy$\\

  \textbf{Covariance, variance d'une somme, corrélation :} 
  La covariance est donnée par : $\mathbb{C}ov(X,Y)=\mathbb{E}(XY)-\mathbb{E}(X)\mathbb{E}(Y)$\\
  (Rappel) : si X et Y sont indé. $\implies$ covariance = 0\\
  Corrélation : \\$\mathbb{C}orr(X,Y)$ ou $\rho(X,Y)=\frac{\mathbb{C}ov(X,Y)}{\sqrt{\mathbb{V}ar(X)\mathbb{V}ar(Y)}}$\\

  \textbf{Espérance et variance conditionnelles :}
  Déf. $\mathbb{E}$:\\
  $\mathbb{E}(X|Y=y)=\sum_{x}xP(X=x|Y=y)$ (cas discret)\\
  $\mathbb{E}(X|Y=y)=\int xf_{X|Y}(x|y)dx$ (cas continu)\\
  avec $\mathbb{E}(\mathbb{E}(X|Y))=\mathbb{E}(X)$\\
  Déf. $\mathbb{V}ar$:\\
  $\mathbb{V}ar(X)=\mathbb{E}(\mathbb{V}ar(X|Y))+\mathbb{V}ar(\mathbb{E}(X|Y))$\\

  \newpage
  
  \textbf{Fcts génératrices des moments et fcts caractéristiques :}
  La FGM est donnée par : $M_X(t)=\mathbb{E}(e^{tX})$\\
  Pour discret : $\sum_{x,p_X(x)>0}e^{tX}P_X(x)$\\
  Pour continu : $\int e^{tX}f_X(x)dx$\\
  (Rmq) : Pour les deux cas, $M_X(0)=1$ et $M^{(p)}(0)=\mathbb{E}(X^p)$.
  Donc,\\
  $M'(0)=\mathbb{E}$ et $\mathbb{V}ar=M''(0)-(M'(0))^2$.\\
  Les fcts caractéristiques sont données par : \\$\phi_X(t)=\mathbb{E}(e^{itX})$ soit $\phi_X(t)=M_X(it)$\\
  Ex. :\\
  \begin{enumerate}
    \item Binomiale : $\phi(t)=(pe^{it}+1-p)^n$
    \item Poisson : $\phi(t)=e^{\lambda(e^{it}-1)}$
    \item Exponentielle : $\phi(t)=(1-\frac{it}{\lambda})^-1$
    \item Normale ($\mu$,$\sigma^2$) : $\phi(t)=e^{i\mu t-\frac{\sigma^2 t^2}{2}}$
  \end{enumerate}


  \section{\underline{Théorèmes limites}}

  \textbf{Inégalité de Markov :}
  Soit X une variable aléa. à valeurs non-négatives. $\forall a>0$, on a :\\
  $\mathbb{P}(X>a)\leq \frac{\mathbb{E}(X)}{a}$\\

  \textbf{Inégalité de Tchecychev :}
  Soit X une variable aléa., de moyenne $\mu$ et de variance $\sigma^2$. $\forall k>0$, on a :\\
  $\mathbb{P}(|X-\mu|>k)\leq \frac{\sigma^2}{k^2}$\\

  \textbf{Inégalité de Jensen :}
  Soit f une fonction convexe, alors :\\
  $\mathbb{E}(f(X))\geq f(\mathbb{E}X)$\\

  \textbf{Autres inégalités :}\\
  Soit X une variable aléa. de moyenne \underline{$\mu=0$} et de variance $\sigma^2$. $\forall a>0$, on a :\\
  $\mathbb{P}(X \geq a)\leq \frac{\sigma^2}{\sigma^2 + a^2}$\\
  Soit X une variable aléa. de moyenne \underline{$\mu$}  et de variance $\sigma^2$. $\forall a>0$, on a :\\
  $\mathbb{P}(X\geq \mu +a)\leq \frac{\sigma^2}{\sigma^2 + a^2}$\\
  et\\
  $\mathbb{P}(X\leq \mu -a)\leq \frac{\sigma^2}{\sigma^2 + a^2}$\\

  \textbf{Loi faible des grands nombres :}
  Pour n variables aléatoires, on suppose que la moyenne $\mu = \mathbb{E}(X_i), i=1,... ,n$ est finie. $\forall\epsilon >0$, on a :\\
  $\mathbb{P}\bigg(\bigg|\frac{\sum_{i=1}^{n}X_i}{n}-\mu\bigg|>\epsilon\bigg)\rightarrow 0$, si $n\rightarrow +\infty$\\
  La convergence vient avec une probabilité.

  \textbf{Loi forte des grands nombres :}
  Pour n variables aléatoires, on suppose que la moyenne $\mu = \mathbb{E}(X_i), i=1,... ,n$ est finie. Avec probabilité égale à 1, on a :\\
  $\frac{\sum_{i=1}^{n}X_i}{n}\rightarrow\mu$, si $n\rightarrow +\infty$\\
  La convergence est sûre.

  \textbf{Théorème central limite (TCL) :}
  Pour n variables aléatoires indé., moyenne $\mu$ et de variance $\sigma^2$. Alors on a :\\
  $\frac{\sum_{i=1}^{n}X_i-n\mu}{\sigma\sqrt{n}}\rightarrow N(0,1)$, si $n\rightarrow +\infty$\\
  On utilise la table de loi normale pour finaliser le tout.
\end{multicols}
\end{document}